\documentclass[11pt]{article}
\usepackage[a4paper,margin=1in]{geometry}
\usepackage{amsmath,amssymb,mathtools,bm}
\usepackage{enumitem,booktabs,array}
\usepackage{tcolorbox}
\usepackage{hyperref}
\usepackage[protrusion=true,expansion=false]{microtype}
\hypersetup{colorlinks=true,linkcolor=blue,urlcolor=blue,citecolor=blue}
\setlist[itemize]{topsep=2pt,itemsep=2pt}
\setlist[enumerate]{topsep=2pt,itemsep=2pt}
\newtcolorbox{keybox}{colback=blue!3!white,colframe=blue!40!black,boxrule=0.4pt,arc=1mm}
\newtcolorbox{alertbox}{colback=red!3!white,colframe=red!40!black,boxrule=0.4pt,arc=1mm}
\newtcolorbox{answerbox}{colback=green!3!white,colframe=green!40!black,boxrule=0.4pt,arc=1mm}
\newtcolorbox{drillbox}{colback=yellow!5!white,colframe=orange!60!black,boxrule=0.4pt,arc=1mm}
\newcommand{\EV}{\mathbb{E}}
\newcommand{\Var}{\mathrm{Var}}
\newcommand{\Cov}{\mathrm{Cov}}
\newcommand{\blank}[1]{\underline{\hspace{#1}}}

\title{E01-01: Coase (1937, Economica) \\
\large ``The Nature of the Firm'' --- Active Recall Workbook}
\author{Empirical Corporate Finance II --- Exam Prep}
\date{\today}
\begin{document}\maketitle

\begin{keybox}
\textbf{Paper in one sentence.} Firms exist because using the price mechanism is costly (search, negotiation, contracting): the firm substitutes internal, authority-based allocation for market transactions wherever the marginal cost of an extra internal transaction is below the marginal cost of a market transaction --- defining firm boundaries at that margin.

\textbf{Placement.} Foundational paper for the theory of the firm. The starting point for the entire transaction-cost literature (Williamson 1979, 1985), incomplete contracts (Grossman-Hart 1986, Hart-Moore 1990), and the organizational-economics branch of corporate finance.
\end{keybox}

\section*{Round 1: Complete Walk-Through}

\subsection*{1.1 The puzzle}
If markets allocate resources efficiently (Smith, Marshall), why do firms exist at all? Why is all production not organized as contracting between individuals through the price system?

\subsection*{1.2 Coase's answer: transaction costs}
Using the price mechanism has costs:
\begin{itemize}
\item Discovering the relevant prices.
\item Negotiating and drafting separate contracts for each transaction.
\item Risk and uncertainty about future contingencies.
\item Tax treatment of market vs.\ internal transactions.
\end{itemize}
Within a firm, these costs are avoided: the entrepreneur \emph{directs} resources by authority, under a single open-ended employment contract.

\subsection*{1.3 Firm boundary condition}
A firm will expand to the point where
\[
\text{marginal cost of organizing one more transaction internally} = \text{marginal cost of doing it through the market}.
\]
Internal cost rises with firm size (decreasing returns to management, errors, mis-allocation); at some point, internal is no longer cheaper than market.

\subsection*{1.4 Distinctive contribution}
\begin{itemize}
\item Firms are not merely production functions; they are \emph{governance structures}.
\item The employment relationship --- an incomplete, open-ended, authority-based contract --- is the defining feature.
\item The price system and the firm are \emph{substitute} coordination mechanisms.
\end{itemize}

\subsection*{1.5 Influence}
\begin{itemize}
\item Williamson (1975, 1985): asset specificity, opportunism, hierarchy --- the transaction-cost branch.
\item Klein-Crawford-Alchian (1978): quasi-rents and holdup --- concretizes Coase's transaction costs.
\item Grossman-Hart (1986), Hart-Moore (1990): property rights over non-contractible actions --- formalizes firm boundary.
\item Jensen-Meckling (1976): agency costs --- complementary lens on internal-vs-market choice.
\end{itemize}

\subsection*{1.6 Caveats}
\begin{alertbox}
\begin{itemize}
\item Paper is prose-theoretic, not formal. The marginal condition is intuition, not a closed model.
\item ``Transaction costs'' are loosely defined; later work tries to pin them down (asset specificity, enforcement).
\item Does not explain \emph{internal organization} of the firm (hierarchy, incentives) --- just firm existence.
\end{itemize}
\end{alertbox}

\section*{Round 2: Level 1 Fill-in}
\begin{enumerate}
\item Core answer to ``why firms?'': \blank{3cm} of using the price mechanism.
\item Firm boundary condition: marginal cost internal $=$ marginal cost \blank{1.5cm}.
\item Firms substitute \blank{1.5cm}-based allocation for \blank{1.5cm}-based allocation.
\item Key contract inside the firm: the \blank{2cm} contract (open-ended, authority).
\item Coase (1937) planted the seed for Williamson and the \blank{2cm}-costs literature.
\end{enumerate}

\section*{Round 3: Level 2 Prompts}
\begin{enumerate}
\item State and explain Coase's transaction-cost argument for firm existence.
\item Derive the boundary condition; what determines when a firm is ``too big''?
\item Contrast Coase with the neoclassical production-function view of the firm.
\item How did Williamson and Klein-Crawford-Alchian sharpen Coase's insight?
\item In an incomplete-contracts framework (Grossman-Hart), how is Coase's boundary reinterpreted?
\end{enumerate}

\section*{Round 4: Minimal Scaffolding}
From a blank page: (i) puzzle, (ii) transaction-cost answer, (iii) boundary condition, (iv) contract structure, (v) descendants.

\section*{Round 5: Blank Page}
Essay: why do firms exist? Coase's answer and its modern refinements.

\vspace{6pt}
\begin{drillbox}
\textbf{Speed Drill.} (1) Main question. (2) Main answer. (3) Boundary condition. (4) Key contract type. (5) Two descendants.
\end{drillbox}

\section*{Quick Reference \& Answer Key}
\begin{answerbox}
\begin{itemize}
\item \textbf{Puzzle.} If markets are efficient, why firms?
\item \textbf{Answer.} Using the price mechanism is costly.
\item \textbf{Boundary.} MC internal $=$ MC market.
\item \textbf{Contract.} Open-ended employment / authority.
\end{itemize}
\end{answerbox}

\begin{keybox}
\textbf{30-second exam pitch.} Coase (1937) asks why production is organized in firms at all if markets work. His answer: using the price mechanism has positive costs --- discovering prices, negotiating contracts, bearing uncertainty, tax treatment. Firms substitute internal, authority-based allocation (directed by an entrepreneur) for market transactions, avoiding these transaction costs. A firm expands until the marginal cost of organizing one more transaction internally equals the marginal cost of the equivalent market transaction --- defining the firm's boundary. The key contract is open-ended employment, under which labor accepts authority within limits in exchange for pay. The paper founds the transaction-cost view of the firm, feeding directly into Williamson (1975, 1985), Klein-Crawford-Alchian (1978), and Grossman-Hart (1986) on property rights. In modern corporate finance, Coase is the grounding for asking: which activities belong inside vs.\ outside firm boundaries, and why?
\end{keybox}

\end{document}
